% ============================================================
% 04-experiments.tex
% ============================================================
\section{Experimental Results}
\label{sec:exp}

This section presents the empirical evaluation of \system{}.

%------------------------------------------------
\subsection{Experimental Benchmark and Protocols}
\label{sec:dataset}

To support empirical evaluation across representative counseling scenarios, our evaluation is established upon a corpus of $200$ administrative queries constructed via a document-grounded semi-automated protocol over official Can Tho University (CTU) statutes. Candidate queries covering diverse conversational phrasing (e.g., student colloquialisms, prerequisite rules, fee schedules) were generated from statutory articles and underwent author-guided review against official decrees to verify ground-truth intents, gold reference passages, and factual answers, with lexical and semantic duplicate checks applied to ensure query novelty. Following standard benchmark practice, this corpus is partitioned into: (1)~\textbf{Development Benchmark ($N=150$ queries):} Spanning 9 institutional categories (\texttt{actual\_tuition}, \texttt{academic\_rules}, \texttt{scholarship}, \texttt{student\_loan}, \texttt{other}, \texttt{social\_support}, \texttt{academic\_program}, \texttt{exemption\_policy}, \texttt{exemption\_basis}), used for offline index calibration, prompt engineering, and qualitative failure analysis; and (2)~\textbf{Held-Out Test Set ($N=50$ queries):} A separated, author-reviewed evaluation suite with exact category counts of 20, 6, 5, 4, 4, 4, 3, 2, and 2, respectively. For domain-stratified analysis, the held-out categories are grouped by assigned specialist into Academic ($N=3$), Financial ($N=24$), Scholarship ($N=5$), and General ($N=18$). All reported primary benchmark metrics in Tables~\ref{tab:retrieval_evaluation} and~\ref{tab:e2e_evaluation} are evaluated on this Held-Out Test Set.

\noindent\textbf{Execution and Evaluation Invariants:} Generative and automated evaluation tasks were executed using Google Gemini 2.5 Flash Lite via Google Cloud Vertex AI under zero-temperature greedy decoding ($T=0.0$). To account for subtle API-level non-determinism across cloud endpoints, each query was evaluated across $R=3$ repeated evaluation runs, yielding $1,050$ end-to-end generation and Ragas evaluation turns ($7\text{ configurations} \times 50\text{ queries} \times 3\text{ runs}$).

\noindent\textbf{Computational Infrastructure:} Neural cross-encoder reranking was executed on an NVIDIA RTX 1060 GPU (4~GB VRAM). Relational curricula are hosted in Neo4j 5 Enterprise with APOC extensions; document parent chunks (2,800 characters, 100 overlap) reside in PostgreSQL, while child chunks (400 characters, 50 overlap) are indexed in Qdrant with \texttt{vietnamese-bi-encoder} (768-d). Neural candidate reranking uses \texttt{bge-reranker-v2-m3}. Multi-agent coordination was orchestrated via LangGraph under a StateGraph supervisor with Redis 7 caching.

%------------------------------------------------
\subsection{Evaluation Metrics}
\label{sec:metrics}

\noindent\textbf{Retrieval and Multi-Agent Benchmarks:} Retrieval effectiveness (Scenario~1) is measured via standard ranking metrics on top-$k$ returned source documents: Hit@$k$ ($k \in \{1, 3\}$), Precision@5, Recall@5, and Mean Reciprocal Rank (MRR@10); domain-stratified analysis reports Hit@1. Multi-agent coordination (Scenario~3) is evaluated across isolated operational decisions: Routing Agent Accuracy and Macro-F1 across the four domains, Intent Classification Accuracy, Tool Selection and Argument Exact Match (EM), End-to-End Pass Rate, Adversarial Tool-Gating Robustness, and Safe Result Behavior under strict operational limits ($\text{max\_tool\_calls}=1, \text{timeout}=60\,\text{s}$).

\noindent\textbf{End-to-End Generative Evaluation:} End-to-end synthesis (Scenario~2) is evaluated using the Ragas framework~\cite{es2023ragas} alongside ground-truth regulatory alignment over Top-7 context: Answer Relevancy (AR), Context Recall (CR), Context Precision (CP), and Answer Correctness (AC). In addition, we evaluate Factual Exact Match (Fact EM) for strict numerical tuition fee and credit agreement, alongside Source Average Precision (Source AP) for verified regulatory citations. Evaluating generative QA through LLM-as-a-judge is well-established~\cite{zheng2023judging,es2023ragas}; to mitigate evaluation bias, our framework enforces reference-grounded metrics conditioned on gold documents, greedy decoding ($T\!=\!0.0$) across $R\!=\!3$ repeated runs, and a programmatic Fact EM evaluator for tuition arithmetic.

%------------------------------------------------
\subsection{Experimental Scenarios and Benchmark Results}
\label{sec:scenarios}

To address the three research questions, we evaluate \system{} across three coordinated experimental scenarios: retrieval component stacking (RQ3), end-to-end QA generation and modular ablation (RQ2), and multi-agent coordination and tool reliability (RQ1).

\subsubsection{Scenario 1 --- Retrieval Component Stacking (E1--E5):}
\emph{Objective \& RQ Alignment:} Scenario~1 directly addresses \textbf{RQ3} by comparing five retrieval configurations at aggregate and domain-stratified levels on the 50-query Held-Out Test Set.
\emph{Configurations:} We compare five cumulative configurations (E1--E5, detailed in Table~\ref{tab:retrieval_evaluation}) transitioning from single sparse/dense baselines to hybrid fusion, neural cross-reranking, and our proposed intent-guided evidence stack.

\emph{Results:} As shown in Table~\ref{tab:retrieval_evaluation}, E5 attains the best Hit@1 ($0.9600$; 48 of 50 queries rank a gold source document first) and MRR@10 ($0.9700$), improving Hit@1 over E1 by 6.0 percentage points. E5 ties E4 at $0.9800$ Hit@3, whereas E1 retains the highest P@5 ($0.2320$) and Recall@5 ($0.9667$). Within the shared hybrid candidate pipeline, adding the cross-encoder from E3 to E4 raises Hit@1 from $0.6800$ to $0.8800$ and MRR@10 from $0.7912$ to $0.9267$. At the domain-stratified level, E5 achieves 1.000 Hit@1 on Academic ($N=3$) and Financial ($N=24$), 0.800 on Scholarship ($N=5$), and 0.944 on General ($N=18$). The E4-to-E5 contrast improves Hit@1 and MRR@10 while leaving Hit@3, P@5, and Recall@5 unchanged; because E5 jointly adds intent governance, graph grounding, and parent--child mapping, this contrast represents the integrated bundle rather than an isolated routing effect.

\begin{table}[!htbp]
\centering
\caption{Scenario 1: Progressive retrieval component stacking on the Held-Out Test Set ($N=50$).}
\label{tab:retrieval_evaluation}
\resizebox{\columnwidth}{!}{
\begin{tabular}{lccccc}
\toprule
\textbf{Retrieval Configuration} & \textbf{Hit@1} & \textbf{Hit@3} & \textbf{P@5} & \textbf{Recall@5} & \textbf{MRR@10} \\
\midrule
E1: BM25 Lexical Baseline        & 0.9000 & 0.9400 & \textbf{0.2320} & \textbf{0.9667} & 0.9333 \\
E2: Dense Vector Semantic        & 0.3800 & 0.6600 & 0.1560 & 0.6433 & 0.5237 \\
E3: Hybrid RRF (BM25 + Dense)    & 0.6800 & 0.9200 & 0.2160 & 0.9133 & 0.7912 \\
E4: Hybrid + Neural Reranker     & 0.8800 & \textbf{0.9800} & 0.2200 & 0.9433 & 0.9267 \\
\textbf{E5: CTU-Chat Proposed Retrieval} & \textbf{0.9600} & \textbf{0.9800} & 0.2200 & 0.9433 & \textbf{0.9700} \\
\bottomrule
\end{tabular}
}
\end{table}

\begin{table}[!htbp]
\centering
\caption{Scenario 2: End-to-end QA generation and modular ablation on Held-Out Test Set ($N=50, R=3$, 1,050 total evaluation turns, Top-7 context).}
\label{tab:e2e_evaluation}
\resizebox{\columnwidth}{!}{
\begin{tabular}{llcccccc}
\toprule
\textbf{Cfg.} & \textbf{Ablation Variant Description} & \textbf{Fact EM} & \textbf{CP} & \textbf{CR} & \textbf{AR} & \textbf{AC} & \textbf{Source AP} \\
\midrule
T1 & BM25 Lexical Baseline                       & 0.5572 & 0.6674 & 0.8800 & 0.6151 & \textbf{0.6735} & 0.8954 \\
T2 & Dense Semantic Baseline                     & 0.5007 & 0.2770 & 0.3400 & 0.2809 & 0.3841 & 0.4651 \\
T3 & Hybrid RRF Baseline                         & 0.5451 & 0.4339 & 0.7733 & 0.5754 & 0.6016 & 0.7490 \\
\textbf{T4} & \textbf{CTU-Chat Full Stack (Proposed)}    & \textbf{0.6203} & \textbf{0.6977} & \textbf{0.8933} & \textbf{0.6195} & 0.6637 & \textbf{0.9311} \\
T5 & w/o Cross Reranker (\texttt{bge-m3})        & 0.5823 & 0.6297 & 0.8133 & 0.5263 & 0.6136 & 0.8168 \\
T6 & w/o Neo4j Knowledge Graph Grounding         & 0.5711 & 0.5683 & 0.8600 & 0.6173 & 0.6582 & 0.9011 \\
T7 & w/o Subspace Governance (Ungoverned)        & 0.6089 & 0.6947 & 0.8600 & 0.5815 & 0.6430 & \textbf{0.9311} \\
\bottomrule
\end{tabular}
}
\smallskip
\parbox{\columnwidth}{\footnotesize\emph{Metrics:} Fact EM: Factual Exact Match on tuition/credits; CP: Context Precision; CR: Context Recall; AR: Answer Relevancy; AC: Answer Correctness; Source AP: Source Average Precision. Normalized to $[0,1]$; higher is better.}
\end{table}

%------------------------------------------------
\subsubsection{Scenario 2 --- End-to-End QA Generation and Modular Ablation (T1--T7):}
\emph{Objective \& RQ Alignment:} Scenario~2 directly addresses \textbf{RQ2} by benchmarking the proposed full system (T4) against three single-agent RAG baselines (T1--T3) and three modular ablations (T5--T7) on the Held-Out Test Set ($1,050$ total evaluation turns; Table~\ref{tab:e2e_evaluation}).
\emph{Results \& Interpretation:} As reported in Table~\ref{tab:e2e_evaluation}, T4 improves on five of six metrics relative to the BM25 baseline, with the largest gain in Fact~EM ($62.03\%$ vs.\ $55.72\%$, $+6.31$ percentage points), while BM25 retains a marginal lead in Answer Correctness. The complementary contributions of modular cross-reranking (T5), graph grounding (T6), and subspace governance (T7) are analyzed in detail under Discussion (Section~\ref{sec:discussion}).

%------------------------------------------------
\subsubsection{Scenario 3 --- Multi-Agent Routing, Tool Reliability, and Robustness:}
\emph{Objective \& RQ Alignment:} Scenario~3 directly addresses \textbf{RQ1} by evaluating multi-agent coordination, tool execution reliability, and bounded safety independently from retrieval or generation.
\emph{Experimental Setup:} As reported in Table~\ref{tab:scenario3_results}, evaluation spans three suites ($N=550$ runs): (1)~\textbf{Supervisor Routing (Panel A):} 100 single-domain benchmark queries ($N=300$ for LLM Supervisor vs. $N=100$ for Rule-Based Router); (2)~\textbf{Specialist Tool Reliability (Panel B):} 30 test cases ($N=90$ runs across 3 repeated runs) covering structured tuition lookup, scholarship calculation, and deterministic tuition reduction arithmetic; and (3)~\textbf{Adversarial Robustness (Panel C):} 20 failure stress tests ($N=60$ runs across 3 repeated runs) challenging boundary invariants.

\begin{table}[!htbp]
\centering
\caption{Scenario 3: Multi-agent routing, specialist tool execution reliability, and robustness under a maximum of one tool call per decision ($\text{max\_tool\_calls}=1$).}
\label{tab:scenario3_results}
\resizebox{\columnwidth}{!}{
\begin{tabular}{lccccccc}
\toprule
\multicolumn{8}{l}{\textbf{Panel A: Supervisor Routing and Specialist Selection ($N=100$ queries; Rule: 1 rep / LLM: 3 reps, $N=300$)}} \\
\midrule
\textbf{Router Variant} & \textbf{Runs ($N$)} & \textbf{Agent Acc.} & \textbf{Macro-F1} & \textbf{Intent Acc.} & \textbf{Bounded} & \multicolumn{2}{c}{\textbf{Routing Paradigm}} \\
\midrule
Rule-Based Router          & 100 & 85.0\% & 69.2\% & 69.0\% & 100.0\% & \multicolumn{2}{c}{Heuristic Regex} \\
LLM Supervisor (\system{}) & 300 & \textbf{97.0\%} & \textbf{97.8\%} & \textbf{95.0\%} & \textbf{100.0\%} & \multicolumn{2}{c}{Structured Pydantic} \\
\midrule
\multicolumn{8}{l}{\textbf{Panel B: Specialist Tool Execution Reliability ($N=90$ runs across 30 test cases, 3 repeated runs)}} \\
\midrule
\textbf{Specialist Tool Function} & \textbf{Runs ($N$)} & \textbf{Selection} & \textbf{Arg. EM} & \textbf{Result Acc.} & \textbf{E2E Pass} & \textbf{Bounded} & \textbf{Paradigm} \\
\midrule
\texttt{tuition\_lookup}                & 30 & 100.0\% & 100.0\% & 100.0\% & 100.0\% & 100.0\% & Structured Catalog \\
\texttt{scholarship\_calculation}       & 30 & 86.7\%  & 76.7\%  & 86.7\%  & 76.7\%  & 100.0\% & LLM Tool Selection \\
\texttt{tuition\_reduction\_calculation}& 30 & 100.0\% & 100.0\% & 100.0\% & 100.0\% & 100.0\% & Arithmetic Tool \\
\midrule
\multicolumn{8}{l}{\textbf{Panel C: Robustness and Bounded Failure Handling ($N=60$ runs under malformed/adversarial inputs)}} \\
\midrule
\textbf{Adversarial Gate} & \textbf{Runs ($N$)} & \textbf{Agent Acc.} & \textbf{Tool Gate} & \textbf{Safe Result} & \textbf{E2E Pass} & \textbf{Bounded} & \textbf{Notes} \\
\midrule
Failure Gate Robustness    & 60 & 100.0\% & 100.0\% & 98.3\% & 98.3\% & 100.0\% & -- \\
\bottomrule
\end{tabular}
}
\smallskip
\parbox{\columnwidth}{\footnotesize\emph{Note:} \texttt{tuition\_lookup} evaluates deterministic structured catalog queries; scholarship and tuition-reduction cases evaluate LLM-driven parameter extraction and external deterministic tool execution.}
\end{table}

\subsection{Discussion}
\label{sec:discussion}

\textbf{Answering RQ1 (Domain-Specialized Multi-Agent Architecture):} Centralizing routing under a structured supervisor with asymmetric privileges and bounds ($\text{max\_tool\_calls}=1$, timeout $=60\,\text{s}$) demonstrated reliable coordination under the evaluated configuration: the supervisor achieves $97.0\%$ agent accuracy, $97.8\%$ Macro-F1, and $95.0\%$ intent accuracy across 300 runs, outperforming rule-based routing ($85.0\%$, $69.2\%$). In specialist tool execution (Panel~B), \texttt{tuition\_lookup} and \texttt{tuition\_reduction\_calculation} achieved $100.0\%$ end-to-end pass rates, while \texttt{scholarship\_calculation} reached $76.7\%$ ($86.7\%$ tool selection), with all three functions maintaining $100.0\%$ bounded completion. Furthermore, adversarial testing (Panel~C) verified $100.0\%$ tool-gating accuracy, $98.3\%$ safe-result behavior, and $100.0\%$ bounded completion. To assess architectural scalability, a paired stress-test compared \system{} against a monolithic single-agent baseline: under the baseline registry (11 tools), both architectures exhibited comparable performance (single-agent $93.3\%$ vs. multi-agent $88.3\%$ E2E, difference not statistically significant); however, when 16 synthetic distractor tools were introduced to simulate institutional expansion (27 tools total), the monolithic agent suffered from decision-space pollution, degrading from $98.3\%$ to $91.7\%$ in tool selection. In contrast, domain-partitioned multi-agent maintained $100.0\%$ tool selection (paired $\Delta = +8.3\%$, 95\% bootstrap CI $[+1.7, +15.0]$, $n=60$). The benefit of domain-specialized multi-agent decomposition became more apparent as the tool registry expanded, because partitioned tool access reduced decision-space interference relative to the evaluated monolithic single-agent baseline.

\textbf{Answering RQ2 (Heterogeneous Knowledge Allocation \& Modular Ablations):} The modular ablations provide evidence that representation-specific knowledge allocation contributes to the evaluated context- and answer-quality metrics. Relative to T4, removing Neo4j (T6) is associated with a Context Precision decrease from $0.6977$ to $0.5683$ ($18.5\%$ relative), together with lower Factual Exact Match ($62.03\%$ to $57.11\%$) and Context Recall ($0.8933$ to $0.8600$). Omitting cross-reranking (T5) is associated with lower Answer Relevancy ($0.6195$ to $0.5263$) and Context Recall ($0.8933$ to $0.8133$), while removing subspace governance (T7) is associated with lower Answer Correctness ($0.6637$ to $0.6430$) and Answer Relevancy ($0.6195$ to $0.5815$). These comparisons support the contribution of the evaluated components within the full configuration but do not establish uniform causal effects across query types.

\textbf{Answering RQ3 (Representation-Aware Retrieval Effectiveness):} The configuration sequence reveals that neural reranking principally improves early-rank ordering rather than evidence coverage: E3-to-E4 gains are concentrated in Hit@1 and MRR@10, while the full stack does not surpass the lexical baseline on P@5 or Recall@5. Domain stratification further shows that E5's additional Hit@1 benefit is concentrated in Financial queries; General queries improve over E1 by E4 and then remain unchanged, while neither Academic nor Scholarship improves over E1. This pattern is consistent with BM25 remaining competitive on exact course codes, cohort identifiers, and administrative terminology. Because E5 changes intent governance, graph grounding, and parent--child mapping together, these results compare complete configurations and do not attribute the observed gains causally to routing alone.

%------------------------------------------------
\subsection{Threats to Validity and Limitations}
\label{sec:limitations}

The evaluation is scoped to Can Tho University; multi-institutional validation remains planned. Domain-stratified strata are imbalanced, so small categories do not support claims of uniform cross-domain improvement. Due to budgetary constraints, Gemini 2.5 Flash Lite served as both generator and judge, introducing potential evaluation bias. Curricula graphs rely on offline schemas, and automated judges penalize valid abstentions. Fault injection across external infrastructure (Neo4j, Qdrant) is reserved for future work.
